\section{Результаты экспериментов}

\subsection{Автоматический режим}
При правильно подобранных коэффициентах, удовлетворяющих условию (7), наблюдается асимптотическая устойчивость:

% \begin{figure}[H]
% 	\centering
% 	\includegraphics[width=0.45\linewidth]{05-angle.png}
% 	\includegraphics[width=0.45\linewidth]{05-position.png}
% 	\caption{Переходные процессы: угол \(\varphi(t)\) (слева) и положение точки опоры \(u(t)\) (справа)}
% 	\label{fig:transient}
% \end{figure}

\subsection{Влияние параметров}
\begin{itemize}
	\item \textbf{Увеличение \(a\)}: ускоряет возврат к вертикали, но при слишком больших значениях могут возникнуть колебания.
	\item \textbf{Увеличение \(b\)}: подавляет колебания, делает процесс апериодическим.
	\item \textbf{Увеличение \(c\)}: ускоряет возврат точки опоры к нулю.
	\item \textbf{Увеличение \(d\)}: демпфирует движение точки опоры.
\end{itemize}

\subsection{Ручной режим}
При ручном управлении удается удерживать маятник вблизи вертикали только при медленном изменении. Одновременная стабилизация точки опоры вручную практически невозможна.